\documentclass[11pt]{article}
\usepackage[margin=1in]{geometry}
\usepackage{float}
\usepackage{array}
\usepackage[hidelinks]{hyperref}
\usepackage{enumitem}
\setlist{itemsep=1pt, topsep=3pt}
\setlength{\parindent}{0pt}
\setlength{\parskip}{7pt}
\setlength{\emergencystretch}{2em}
\renewcommand{\arraystretch}{1.25}
\newcommand{\code}[1]{\texttt{#1}}
\newcolumntype{P}[1]{>{\raggedright\arraybackslash}p{#1}}

\begin{document}

\begin{center}
{\Large DS-GA 1011 Assignment 4}\\[4pt]
Natiq Khan (nak9135) \quad Fall 2026
\end{center}

\textbf{Code:} \url{https://github.com/nak9135/nlp_a4}\\
\textbf{Model checkpoints:}\\ \url{https://drive.google.com/drive/folders/1j2nShBflPWAMhAdkVlnazR8XTUFrPFTk}\\
(\code{t5\_ft\_best} is the fine-tuned model for Q7, \code{t5\_scratch\_best} is the extra credit model.)

I trained everything on a free Kaggle notebook with a T4 GPU. I started on the NYU HPC cluster, but
I couldn't get GPU jobs to start there, so I switched to Kaggle to save time. The repo has scripts
for both.

%==========================================================
\section*{Part 1: BERT sentiment classification}

\subsection*{Q1. Fine-tuning BERT}

I wrote the training loop: for each batch, run the model, compute the loss, backpropagate, and
update the weights and learning rate. I used the default settings (3 epochs, batch size 8,
learning rate $5\times10^{-5}$).

\begin{itemize}
  \item Small debug run: 88.3\% accuracy
  \item Full run: \textbf{92.81\%} accuracy
\end{itemize}

\subsection*{Q2a. My transformation: a ``careless reviewer''}

The idea is to make each review look like it was typed quickly by a real person: a few words are
swapped for synonyms and a few have typos. The meaning stays the same.

How it works, word by word:
\begin{enumerate}
  \item Split the review at its \code{<br />} line breaks, then split each part into words with
        NLTK's \code{word\_tokenize}.
  \item Only change a word if it is made of letters, is at least 3 letters long, and is not a
        negation (\emph{not, no, never, \dots}) or a very common word (\emph{the, and, was, \dots}).
        Negations are left alone so the sentiment can't flip.
  \item For each word that can be changed, roll a random number between 0 and 1:
  \begin{itemize}
    \item below 0.25: swap in a \textbf{synonym}, picked at random from WordNet (from the word's
          first two meanings, one-word synonyms only). If there is none, keep the word.
    \item between 0.25 and 0.45: add \textbf{one typo} at a random position (never the first
          letter). The typo is one of: hit a neighbouring keyboard key, swap two letters, drop a
          letter, or double a letter (each equally likely).
    \item otherwise: leave the word alone.
  \end{itemize}
  \item Keep the original capitalization and put the words back together into a sentence with
        NLTK's detokenizer.
\end{enumerate}
The random seed is fixed (0), so the transformed test set is always the same.

\textbf{Example} (real output):\\
\emph{Before:} This movie was absolutely wonderful. The acting was great and I loved the ending,
but the middle dragged a little.\\
\emph{After:} This movie was absolutely wonderfu. The actig was graet and I lovrd the endkng, but
the middle draggd a little.

\textbf{Why this is reasonable:} people really do write reviews like this, especially on their
phones, and different people use different words for the same opinion. A human can still tell
exactly what the reviewer meant, so the label doesn't change. But BERT has a harder time: a typo
like \emph{graet} gets split into word pieces it rarely saw during training, so an important word
like \emph{great} can lose its meaning for the model.

\subsection*{Q2b. Results}

\begin{center}
\begin{tabular}{|l|c|}
\hline
Test set & Accuracy \\ \hline
Original & 92.81\% \\
Transformed & 88.06\% \\ \hline
Drop & \textbf{4.75 points} \\ \hline
\end{tabular}
\end{center}

%==========================================================
\subsection*{Q3. Data augmentation}

I picked 5,000 random training reviews, applied my transformation to them, and added them to the
25,000 normal training reviews (30,000 in total). Then I trained a new BERT model on that with the
same settings as Q1.

\begin{center}
\begin{tabular}{|l|c|c|}
\hline
Model & Original test & Transformed test \\ \hline
Trained on normal data (Q1) & 92.81\% & 88.06\% \\
Trained on augmented data (Q3) & 92.90\% & 90.12\% \\ \hline
Change & +0.09 & \textbf{+2.06} \\ \hline
\end{tabular}
\end{center}

\textbf{Did accuracy on the transformed test set improve?} Yes, from 88.06\% to 90.12\%. That wins
back about 2 of the 4.75 points we lost, but not all of them.

\textbf{What happened on the original test set?} Basically nothing (92.81\% to 92.90\%). So the
augmentation didn't hurt performance on clean reviews.

\textbf{Why it works:} the first model only ever saw clean reviews, so it learned to rely on exact
words like \emph{great} or \emph{terrible}. Once those are misspelled, it gets confused. The
augmented model has seen misspelled and reworded reviews during training, with the right labels,
so it learns that \emph{graet} still means \emph{great}. It doesn't recover fully because only 1
in 6 training reviews was changed, and the test reviews have brand-new random typos.

\textbf{A limitation:} augmentation only helps with the kind of change you train on. Here the
training typos and the test typos come from the same code, which is the easiest possible case.
Real-world shifts (slang, sarcasm, other kinds of reviews) are hard to predict, and this model has
no reason to handle them any better.

\emph{Side note:} my first augmented training run got stuck guessing the same label for every
review (50\% accuracy). This sometimes happens when fine-tuning BERT. Rerunning with a different
random seed fixed it, and that run gave the numbers above.

%==========================================================
\section*{Part 2: T5 text-to-SQL}

\subsection*{Q4. Data statistics}

All counts use the T5 tokenizer. ``Vocabulary'' means how many different tokens show up. The test
set has 432 examples.

\begin{table}[H]
\centering
\begin{tabular}{|l|c|c|c|c|}
\hline
 & \multicolumn{2}{c|}{Table 1: before processing} & \multicolumn{2}{c|}{Table 2: after processing} \\ \hline
 & Train & Dev & Train & Dev \\ \hline
Number of examples & 4,225 & 466 & 4,225 & 466 \\
Average question length & 17.10 & 17.07 & 23.10 & 23.07 \\
Average SQL length & 216.37 & 210.05 & 217.37 & 211.05 \\
Vocabulary (questions) & 791 & 465 & 796 & 470 \\
Vocabulary (SQL) & 555 & 395 & 557 & 397 \\
Longest question & 59 & 43 & 65 & 49 \\
Longest SQL & 510 & 502 & 511 & 503 \\
Unknown tokens in SQL & 510 & 55 & 0 & 0 \\ \hline
\end{tabular}
\end{table}

``After processing'' means I added the prefix \code{translate English to SQL:} to each question,
added an end token to both sides, and replaced the \code{<} sign in SQL (more on that in Q5). That
last change is why the unknown tokens drop to 0.

\subsection*{Q5. How I set up T5 (Table 3)}

\begin{table}[H]
\centering
\begin{tabular}{|P{3.45cm}|P{11.45cm}|}
\hline
\textbf{Data processing} &
I added the prefix \code{translate English to SQL:} to each question, since that's the style T5
was trained with. T5's vocabulary has no \code{<} sign, so the model could never write it. I
replace every \code{<} with the word \code{lessthan} before training and change it back after the
model writes its answer. I checked that every query survives this round trip unchanged. \\ \hline
\textbf{Tokenization} &
The standard \code{t5-small} tokenizer for both questions and SQL. Questions are capped at 128
tokens and SQL at 512 (nothing is long enough to get cut). Each SQL query ends with an end token
so the model learns when to stop. During training the decoder is fed the correct SQL shifted by
one token, starting from T5's usual start token. Padding is ignored in the loss. \\ \hline
\textbf{Architecture} &
The pretrained \code{t5-small} model, with all of its weights trained (nothing frozen). \\ \hline
\textbf{Hyperparameters} &
AdamW optimizer, learning rate $5\times10^{-4}$, weight decay 0.01, batch size 16. The learning
rate warms up for 1 epoch and then slowly decreases (cosine schedule). Up to 20 epochs. After each
epoch I generate SQL for the dev set (greedy decoding), run it on the database, and keep the
version with the best F1. Training would stop early after 5 epochs without improvement, but that
didn't happen. The best version came from epoch 17. Training took about 2 hours. \\ \hline
\end{tabular}
\end{table}

\subsection*{Q6. Results (Table 4)}

\begin{table}[H]
\centering
\begin{tabular}{|l|c|c|c|}
\hline
 & Query EM & Record EM & F1 \\ \hline
\multicolumn{4}{|l|}{\textbf{Dev set}} \\
Final model (best epoch, 17 of 20) & 68.45 & 82.83 & \textbf{84.64} \\
Same run after 5 epochs & 47.21 & 64.16 & 70.22 \\
Same run after 10 epochs & 63.73 & 77.25 & 79.07 \\
First run (before my timeout fix) & 65.88 & 80.04 & 81.73 \\ \hline
\multicolumn{4}{|l|}{\textbf{Test set}} \\
Final model & n/a* & 77.31 & \textbf{79.47} \\ \hline
\end{tabular}
\end{table}
{\small *The Gradescope autograder shows 0.00 for test Query EM, which looks like it simply
isn't computed for the test set, so I left it as n/a.}

I had limited GPU time, so I trained one setup instead of comparing many. The middle rows show how
the score improved over the epochs. My first run used the same settings, but it had a bug in
evaluation: SQL queries that ran too long were not actually stopped, which slowed everything down
and made correct answers count as timed out. That made the score look like it dropped and stopped
training early. After I fixed it (any query is now stopped after 60 seconds), the same setup
trained for all 20 epochs.

Query EM is lower than Record EM because 67 dev queries are written differently from the correct
query but still return the same results.

\subsection*{Q6. Error analysis (Table 5)}

80 of the 466 dev queries return the wrong results. Here is what goes wrong:

\begin{table}[H]
\centering\small
\begin{tabular}{|P{2.4cm}|P{4.6cm}|P{5.5cm}|P{2.2cm}|}
\hline
\textbf{Error type} & \textbf{Example} & \textbf{What goes wrong} & \textbf{How often} \\ \hline
Wrong value &
``us air flights departing from charlotte around 1pm'': the model writes 1330--1330 instead of
1230--1330 &
The query has the right shape but a time, city, code or date is wrong. Vague phrases like
``around 1pm'' and rare airport codes are the hardest. &
45 of 466 \\ \hline
Missing or extra parenthesis &
``show me the latest flight to love field'': correct query, but with one extra \code{)} at the
end &
The correct queries have lots of nested brackets, and the model loses count, so the query can't
run. &
13 of 466 \\ \hline
Made-up table names &
``ground transportation in philadelphia in the morning on wednesday'': uses a table name it never
defined &
The model copies a pattern it has seen a lot but forgets to list one of the tables, or lists one
twice. A few queries keep repeating conditions until they time out. &
19 of 466 (14 name errors, 5 timeouts) \\ \hline
Wrong columns &
``define airline us'': returns only the airline code instead of code, name and note &
The conditions are right, but the query returns the wrong columns or is built differently.
These questions are worded in unusual ways. &
10 of 466 \\ \hline
Bad labels in the data &
``list all flights from denver to philadelphia two saturdays from now'': the ``correct'' answer is
a dummy query about a state code ``ASD'' &
Some correct answers in the data are placeholders unrelated to the question. The model's answer
makes sense but is marked wrong. &
5 of 466 \\ \hline
\end{tabular}
\end{table}

(In total, 32 dev queries fail to run; 9 of those still count as correct because the right answer
is empty.)

%==========================================================
\section*{Extra credit: T5 trained from scratch}

Same \code{t5-small} model, but starting from random weights instead of pretrained ones. Without
pretraining the model has to learn everything from 4,225 examples, so I made the SQL easier to
learn.

\textbf{My main change: a SQL vocabulary.} The normal T5 tokenizer cuts SQL into lots of tiny
pieces (\code{flight\_1.flight\_id} becomes 8 tokens), so a query is about 217 tokens long. I gave
every word that appears in the training SQL its own token instead (644 words, like \code{SELECT},
\code{flight\_1.flight\_id} or \code{'BOSTON'}). Now a query is only about 61 tokens. If the model
needs a word it never saw in training (like a flight number from the question), it can still spell
it out with normal T5 pieces. The questions still use the normal T5 tokenizer.

\subsection*{Data statistics (new Tables 1 and 2)}

Question lengths use the T5 tokenizer; SQL lengths use my SQL vocabulary.

\begin{table}[H]
\centering
\begin{tabular}{|l|c|c|c|c|}
\hline
 & \multicolumn{2}{c|}{Before processing} & \multicolumn{2}{c|}{After processing} \\ \hline
 & Train & Dev & Train & Dev \\ \hline
Number of examples & 4,225 & 466 & 4,225 & 466 \\
Average question length & 17.10 & 17.07 & 18.10 & 18.07 \\
Average SQL length & 60.90 & 58.90 & 61.90 & 59.96 \\
Vocabulary (questions) & 791 & 465 & 792 & 466 \\
Vocabulary (SQL) & 644 & 393 & 645 & 402 \\
Longest question & 59 & 43 & 60 & 44 \\
Longest SQL & 158 & 146 & 159 & 147 \\
Unknown tokens in SQL & 0 & 0 & 0 & 0 \\ \hline
\end{tabular}
\end{table}

Processing here just adds an end token (no prefix this time). 17 dev queries contain words that
never appear in the training SQL, which is why the dev SQL vocabulary goes up a little.

\subsection*{Model setup (new Table 3)}

\begin{table}[H]
\centering
\begin{tabular}{|P{3.45cm}|P{11.45cm}|}
\hline
\textbf{Data processing} &
No prefix (a model trained from scratch has never seen it). SQL is split on spaces into words.
Known words get their own token; unknown words are spelled with T5 pieces. \\ \hline
\textbf{Tokenization} &
Questions: the normal \code{t5-small} tokenizer, capped at 128 tokens. SQL: my SQL vocabulary
(644 new tokens) plus T5 pieces for unknown words, ending with an end token. Same decoder setup as
in Q5. \\ \hline
\textbf{Architecture} &
The \code{t5-small} architecture with random starting weights. The only change is 644 extra rows
in the embedding table (one per new SQL word). All weights are trained. \\ \hline
\textbf{Hyperparameters} &
AdamW, learning rate $1\times10^{-4}$, weight decay 0.01, batch size 16, gradient clipping at 1.0.
Warm-up for 2 epochs, then a cosine schedule, up to 80 epochs. Every 5 epochs I score the dev set
and keep the best version; training stops after 4 checks in a row with no improvement. It stopped
at epoch 65, and the best version was from epoch 45. Training took about 1.3 hours. \\ \hline
\end{tabular}
\end{table}

\subsection*{Results (new Table 4)}

\begin{table}[H]
\centering
\begin{tabular}{|l|c|c|c|}
\hline
 & Query EM & Record EM & F1 \\ \hline
\multicolumn{4}{|l|}{\textbf{Dev set}} \\
Final model (learning rate $10^{-4}$, with clipping) & 56.22 & 67.60 & \textbf{69.49} \\
Same run after 10 epochs & 38.41 & 48.07 & 51.03 \\
Same run after 25 epochs & 46.57 & 59.44 & 61.16 \\
Learning rate $3\times10^{-4}$, with clipping & 29.61 & 39.91 & 43.54 \\
Learning rate $5\times10^{-4}$, no clipping (failed) & 0.00 & 11.80 & 11.80 \\
Learning rate $10^{-3}$, no clipping (failed) & 0.00 & 11.80 & 11.80 \\ \hline
\multicolumn{4}{|l|}{\textbf{Test set}} \\
Final model & n/a* & 60.42 & \textbf{62.96} \\ \hline
\end{tabular}
\end{table}
{\small *Same as before: the autograder doesn't report Query EM on the test set.}

\textbf{What I tried:} my first two runs used higher learning rates without gradient clipping
(the setup that worked for fine-tuning). They started out fine, but once the learning rate got to
its peak, training fell apart and the models wrote almost the same query for every question.
Adding gradient clipping (which limits how big each update can be) fixed this, and the lower
learning rate worked best. The final model gets 69.49 F1 on dev and 62.96 on test. That's well
below the fine-tuned model (84.64 and 79.47), but above the 50 F1 target.

\subsection*{Error analysis (new Table 5)}

151 of the 466 dev queries return the wrong results.

\begin{table}[H]
\centering\small
\begin{tabular}{|P{2.4cm}|P{4.6cm}|P{5.5cm}|P{2.2cm}|}
\hline
\textbf{Error type} & \textbf{Example} & \textbf{What goes wrong} & \textbf{How often} \\ \hline
From and to swapped &
``show me the flights from dallas to boston'': the model writes a query for flights from Boston to
Dallas &
The biggest problem by far. Everything else in the query is right, but the cities are in the wrong
places. The model knows which cities are mentioned, but not which one follows ``from''. The
fine-tuned model only did this once, so pretraining seems to teach this. &
78 of 466 (about half of all errors) \\ \hline
Wrong value &
``how many seats in a 734'': the model writes aircraft \code{'100'} instead of \code{'734'} &
A city, code, time or number is wrong. Rare codes are the hardest because the model has to copy
them from the question. &
48 of 466 \\ \hline
Broken SQL &
``flights from la guardia or jfk to cleveland'': the query is missing part of the ``or'' and a
closing bracket &
11 queries have unbalanced brackets and 3 use a table name wrongly, so they can't run. This
happens less than with the fine-tuned model (32 queries), probably because each SQL word is a
single token. &
14 of 466 (9 counted as wrong) \\ \hline
Wrong columns or conditions &
``define airline us'': returns only the code instead of code, name and note &
The query returns the wrong columns, or adds, drops or misplaces a condition. &
13 of 466 \\ \hline
Bad labels in the data &
``list the morning flights at a 124 dollars from atlanta to boston'': the ``correct'' answer is the
dummy ``ASD'' query &
Same placeholder answers as in the Q6 table. The model's answer makes sense but is marked wrong. &
3 of 466 \\ \hline
\end{tabular}
\end{table}

\end{document}
